\documentclass{article}%
\usepackage[T1]{fontenc}%
\usepackage[utf8]{inputenc}%
\usepackage{lmodern}%
\usepackage{textcomp}%
\usepackage{lastpage}%
\usepackage{geometry}%
\geometry{margin=2.2cm}%
%
\usepackage{amsmath}%
\usepackage{amssymb}%
\usepackage{booktabs}%
\usepackage{longtable}%
\usepackage{array}%
\usepackage{hyperref}%
\hypersetup{hidelinks}%
\title{IME 2027 {-} Química (claude\_sonnet\_5\_5) x Poliedro\\[4pt]\large Comparação com o gabarito Poliedro Resolve\thanks{Gerado com gpt-resolve: \url{https://github.com/lgabs/gpt-resolve}}}%
\date{09/11/2026}%
%
\begin{document}%
\normalsize%
\maketitle%
\section*{Resumo}\textbf{Nota média: 9,8/10} (22 de 23 itens com a resolução correta; cada item vale o mesmo dentro da questão).\par\medskip%
\begin{longtable}{@{}c c p{0.78\textwidth}@{}}\toprule Questão & Nota & Resumo \\ \midrule
1 & 10,0 & A resolução está correta e coincide com o gabarito do Poliedro: as mesmas duas reações balanceadas (MnO2 + 4 HCl $\to$ MnCl2 + Cl2 + 2 H2O e Cl2 + 2 NaBr $\to$ 2 NaCl + Br2) e m(MnO2) $\approx$ 0,87 g. Recalculei com as massas atômicas da prova e obtive 0,8703 g. Não há divergência a resolver. \\
2 & 10,0 & A resolução da questão 2 está correta e confere com o gabarito Poliedro nos dois itens. Item a: é o mesmo ciclo, isoterma T1, adiabática, isoterma T3 > T1, adiabática, na ordem 1$\to$2$\to$3$\to$4$\to$1; II e IV estão tracejadas por serem irreversíveis, o que é um refinamento. Item b: a demonstração é a mesma, W2 = $\Delta$U2 = -W4 porque U só depende de T, T2 = T1 e T4 = T3. A convenção de sinal (trabalho sobre o gás) é a mesma do gabarito, e as expressões de W1 e W3 são equivalentes. Conferi em Python as coordenadas do esboço e a verificação numérica (W2 = +0,40RT1, W4 = -0,40RT1). Não achei erro. \\
3 & 10,0 & A resolução está certa e bate com o Poliedro nos quatro itens:\newline%
{-} **a)** ordem 2 em SO2 e ordem 0 em O3.\newline%
{-} **b)** v = k{[}SO2{]}$^2$.\newline%
{-} **c)** k = 1,888 $\approx$ 1,89 L$\cdot$mol$^{-}$$^{1}$$\cdot$s$^{-}$$^{1}$.\newline%
{-} **d)** Ea = R$\cdot$ln 4/(7/3000) $\approx$ 594R $\approx$ 4,93 kJ/mol.\newline%
\newline%
Refiz as contas e confirmei: a razão k2/k1 dá 4,00, e só há diferenças de arredondamento. Nenhum erro foi encontrado nem na resolução nem no gabarito. As únicas ressalvas são hipóteses implícitas de baixa importância: a unidade de k a 1000 K e o uso de Arrhenius numa reação sob luz UV. \\
4 & 7,5 & Os itens a, b e d batem com o gabarito Poliedro. No item d, refiz a conta: i = 0,58, W = 0,1156 mol/kg, $\Delta$T = 0,343 K e Tf = 5,11 $^\circ$C. A única divergência é o item c. A resolução dá Kc $\approx$ 164 L/mol, supondo que a solução 0,10 mol/L do item d é o próprio equilíbrio à mesma temperatura. O gabarito afirma que não há valor numérico e responde Kc = 16,4$\cdot$V. As duas respostas são coerentes entre si, porque 164 é exatamente 16,4$\cdot$V com V = 10 L por mol de ácido, e a resolução também apresenta a forma geral 16,4/c0. A diferença vem da falta de concentração no enunciado e do critério de correção, não de erro de conta; por isso classifiquei o item c como 'unclear'. \\
5 & 10,0 & A resolução e o gabarito do Poliedro concordam: a substância é o n{-}heptano. A diferença entre os dois calorímetros vale C($\Delta$T\_V - $\Delta$T\_P) = $\Delta$n$\cdot$RT, o que dá $\Delta$n = 450$\cdot$16,6/(8,3$\cdot$300) = +3. Pela relação $\Delta$n = y/4 - 1, só um hidrocarboneto com 16 H (C7H16) atende. Conferi as contas e os sinais em Python e está tudo certo. A resolução também justifica a exclusão dos outros candidatos: o acetileno tem $\Delta$n negativo, e o propeno e o p{-}xileno dão $\Delta$T\_V - $\Delta$T\_P = 2,8 K e 8,3 K. Ficam só observações menores sobre a idealização do enunciado. \\
6 & 10,0 & Os três itens da resolução batem com o gabarito Poliedro e estão certos. Em (a), a equação é a esterificação 1:1:1:1, com produtos acetato de etila e água. Em (b), $\alpha$ do ácido é 1,6/4,0 = 0,40. Em (c), Kc = (1,6$\cdot$1,6)/(2,4$\cdot$0,4) = 8/3 $\approx$ 2,67, e o volume se cancela porque $\Delta$n = 0. Conferi o enunciado no PDF original e refiz as contas em python, inclusive a quadrática de consistência e as massas molares exatas. Ficam só pontos menores: a água entra em Kc e as massas molares foram arredondadas. \\
7 & 10,0 & A resolução está correta e coincide com o Poliedro nos quatro itens: $\alpha$ = 1, $\beta$ = 2, $\gamma$ = -1; k = 1,25$\times$10$^{-}$$^2$ L$\cdot$mol$^{-}$$^{1}$$\cdot$s$^{-}$$^{1}$; $\nu$ = k{[}A{]}{[}B{]}$^2$/{[}C{]}; ordem global 2. Os dois partem da mesma hipótese de erro de impressão no $\nu$$_{0}$ do exp. 4 (1,28$\times$10$^{-}$$^3$ em vez do impresso 1,28$\times$10$^{-}$$^{6}$). O valor impresso está mesmo no PDF e, usado literalmente, gera expoentes não inteiros. Conferi em Python o determinante (20), a solução do sistema, o valor de k e as velocidades previstas. A dedução por estado estacionário que a resolução apresenta confirma o resultado de forma independente. O único ponto a levar ao professor é como a banca vai tratar o dado inconsistente do enunciado. \\
8 & 10,0 & A resolução está correta e coincide com o gabarito Poliedro nos dois itens. Na reação de Wurtz, Y = R--R é simétrico. O menor Y com dois carbonos quirais e opticamente inativo é o meso{-}3,4{-}dimetil{-}hexano (C8), que vem do sec{-}butila. Logo X é o 2{-}clorobutano. As configurações (3R,4S) dos desenhos em zigue{-}zague e em Fischer foram conferidas numericamente pelas coordenadas 3D, e o confôrmero desenhado tem centro de inversão, o que confirma a forma meso. Ficam só observações de baixa relevância: o rendimento real do Wurtz com haleto secundário, o uso do termo \textbackslash{}"carbono quiral\textbackslash{}" em composto meso e o fato de a imagem do gabarito não estar disponível. Resultado: 2/2 itens em concordância. \\
9 & 10,0 & A resolução está correta e bate com o gabarito do Poliedro: M = 106 g/mol e A = p{-}xileno (1,4{-}dimetilbenzeno). Conferi a crioscopia em python: W = 0,25 mol/kg, n = 0,50 mol, M = 106 g/mol. A eliminação também está certa: o Tollens negativo exclui o benzaldeído, e só o p{-}xileno dá um único monoderivado halogenado no anel e na cadeia lateral. A transcrição confere com a página 5 da prova. Não há erro. As dúvidas são só sobre a leitura de \textbackslash{}"único produto possível\textbackslash{}" (tipo de halogenação, só monoderivados) e sobre hipóteses implícitas razoáveis. \\
10 & 10,0 & Questão 10 (item único): a resolução está certa e equivalente ao gabarito do Poliedro Resolve. As duas usam e\^{}\{$\lambda$t\} = 1,024 com ln 1,024 $\approx$ 0,024, o que dá t/t$\frac12$ = 0,0346 e t $\approx$ 4,88$\times$10$^{8}$ anos. A resolução arredonda para 4,9$\times$10$^{8}$ anos e cita também o valor com o ln exato ($\approx$4,8$\times$10$^{8}$ anos). Não há erro de conteúdo em nenhum dos dois. As diferenças são só de arredondamento, e o enunciado no PDF confere com a transcrição. \\
\bottomrule\end{longtable}%
\section*{Principais pontos de dúvida}%
\begin{itemize}%
\item%
\textbf{Q4} \textit{(statement)} Item c: o enunciado não dá a concentração (ou o volume) do equilíbrio de referência, e Kc = 16,4/c0 depende dela. O gabarito Poliedro deixa a resposta como Kc = 16,4$\cdot$V. A resolução fixa c0 = 0,10 mol/L supondo que a solução do item d é o próprio equilíbrio, e obtém 164 L/mol. Para isso precisa supor a mesma temperatura, o que o enunciado não garante: a solução do item d está no ponto de fusão, $\approx$ 5 $^\circ$C, e a T do equilíbrio original não é dada. O gabarito oficial do IME pode esperar 164 (c0 = 0,10), 16,4 (base 1 L) ou a expressão literal. Vale levar a um professor que critério de correção será aceito.%
\item%
\textbf{Q7} \textit{(statement)} Pelo PDF original, a tabela impressa traz $\nu$$_{0}$ = 1,28$\times$10$^{-}$$^{6}$ mol L$^{-}$$^{1}$ s$^{-}$$^{1}$ no exp. 4. Com esse valor os expoentes saem não inteiros ($\alpha$ $\approx$ -1,49; $\beta$ $\approx$ 1,00; $\gamma$ $\approx$ 0,49). Só com 1,28$\times$10$^{-}$$^3$ (razão 2$^{1}$$^{0}$ = 1024) se chega a (1, 2, -1). A resolução e o gabarito fazem a mesma correção. Convém o professor verificar se a banca publicou errata ou se vai aceitar outras abordagens: usar o valor literal, ou usar só os exp. 1 a 3 com o mecanismo, que já bastam para chegar a (1, 2, -1) se a ordem em A for 1.%
\end{itemize}%
\section*{Questão 1}%
\href{https://poliedroresolve.sistemapoliedro.com.br/vestibulares/ime/2027/ime-2-fase-dia-3-28-10-2026-2-fase-dia-3/dissertativa-1-quimica-2-fase-dia-3-ime-2027}{Resolução do gabarito}\par\smallskip%
\begin{longtable}{@{}p{0.06\textwidth} p{0.33\textwidth} p{0.33\textwidth} p{0.1\textwidth} p{0.08\textwidth}@{}}\toprule Item & Resolução & Gabarito & Veredito & Certo \\ \midrule
unico & MnO2 + 4 HCl $\to$ MnCl2 + Cl2 + 2 H2O; Cl2 + 2 NaBr $\to$ 2 NaCl + Br2; m(MnO2) $\approx$ 0,87 g (0,0100 mol) & MnO2 + 4HCl $\to$ MnCl2 + Cl2 + 2H2O; Cl2 + 2NaBr $\to$ 2NaCl + Br2; m(MnO2) = 0,87 g & igual & ambos \\
\bottomrule\end{longtable}%
\begin{itemize}%
\item%
\textbf{unico)} As duas reações são as mesmas e estão balanceadas, e a massa coincide. Refiz a conta com as massas atômicas da prova: M(NaBr) = 22,990 + 79,904 = 102,894 g/mol, então n(NaBr) = 2,06/102,894 = 0,02002 mol. Daí n(Cl2) = n(MnO2) = 0,01001 mol. Com M(MnO2) = 86,936 g/mol, a massa é 0,8703 g, ou seja, cerca de 0,87 g. O enunciado transcrito confere com o original (página 2).%
\end{itemize}%
\paragraph{Pontos de dúvida}%
\begin{itemize}%
\item%
\textit{(solution)} A resolução supõe, sem deixar isso como premissa central, rendimento de 100\%, HCl em excesso e que a massa pedida é a de MnO2 puro, e não a do minério pirolusita. Esta é a leitura natural do enunciado, que fala em 'massa de óxido de manganês IV'. A própria resolução registra essas hipóteses na conferência.%
\end{itemize}%
\section*{Questão 2}%
\href{https://poliedroresolve.sistemapoliedro.com.br/vestibulares/ime/2027/ime-2-fase-dia-3-28-10-2026-2-fase-dia-3/dissertativa-2-quimica-2-fase-dia-3-ime-2027}{Resolução do gabarito}\par\smallskip%
\begin{longtable}{@{}p{0.06\textwidth} p{0.33\textwidth} p{0.33\textwidth} p{0.1\textwidth} p{0.08\textwidth}@{}}\toprule Item & Resolução & Gabarito & Veredito & Certo \\ \midrule
a & Diagrama P$\times$V do ciclo 1$\to$2$\to$3$\to$4$\to$1: etapa I na isoterma T1 (compressão de V1 a V2), etapa II rápida e adiabática de 2 a 3, com T3 > T1, etapa III na isoterma T3 (expansão de V3 a V4) e etapa IV rápida e adiabática de 4 a 1. As etapas II e IV aparecem tracejadas porque são irreversíveis, mas foram desenhadas mais inclinadas que as isotermas. O gráfico tem setas de sentido e rótulos de estados e etapas. & Gráfico P$\times$V do ciclo 1$\to$2 (I, isoterma T1), 2$\to$3 (II, adiabática), 3$\to$4 (III, isoterma T3), 4$\to$1 (IV, adiabática). & igual & ambos \\
b & Em II e IV, Q = 0, então W2 = $\Delta$U2 = U(T3) - U(T1) e W4 = $\Delta$U4 = U(T1) - U(T3). Com Cv constante, W2 = Cv(T3 - T1) = -W4. Logo W2 + W4 = 0 e W\_tot = W1 + W3, com W1 = RT1 ln(V1/V2) > 0 e W3 = -RT3 ln(V4/V3) < 0 (trabalho realizado sobre o gás). & W2 = n$\cdot$Cv$\cdot$(T3-T1) e W4 = n$\cdot$Cv$\cdot$(T1-T3) = -W2 $\Rightarrow$ Wtot = W1 + W3, com W1 = -nRT1$\cdot$ln(V2/V1) e W3 = -nRT3$\cdot$ln(V4/V3). & equivalente & ambos \\
\bottomrule\end{longtable}%
\begin{itemize}%
\item%
\textbf{a)} O ciclo é o mesmo: duas isotermas (T1 abaixo, T3 acima) ligadas por dois trechos adiabáticos, na mesma ordem e no mesmo sentido. Conferi as coordenadas do TikZ: o estado 3 (0,26; 5,77) e o estado 4 (0,48; 3,125) estão na isoterma 1,5/V. As curvas tracejadas ligam de fato 2$\to$3 e 4$\to$1 (deram 5,766 e 1,000). As setas seguem o sentido certo (I e II para a esquerda, III e IV para a direita). Também estão corretas as relações P2 > P1, P3 > P2, P4 < P3 e P4 > P1, esta última porque P4V4 = RT3 > RT1 = P1V1 e V4 < V1. A única diferença é que a resolução traça II e IV tracejadas, sem usar PV\^{}$\gamma$, por serem rápidas e irreversíveis. Isso é um refinamento e não contradiz o gabarito.%
\item%
\textbf{b)} A demonstração é a mesma: a primeira lei com Q = 0 dá W = $\Delta$U, e U só depende de T. Como T2 = T1 e T4 = T3, W2 e W4 se anulam. A convenção de sinal também é a mesma nos dois (W é o trabalho sobre o gás, $\Delta$U = Q + W): o gabarito escreve W2 = nCv(T3 - T1) > 0 na compressão. Com n = 1, RT1 ln(V1/V2) = -RT1 ln(V2/V1), então o W1 coincide, e o W3 é idêntico. A resolução é um pouco mais geral porque usa U(T) sem precisar de Cv constante. Refiz a verificação numérica (gás monoatômico, V2 = V1/3, P\_ext = 5P1 em II e P1 em IV) com frações exatas: T3 = 19T1/15, V3 = 19V1/75, V4 = 0,6V1, W2 = +0,40RT1 e W4 = -0,40RT1. Tudo confere.%
\end{itemize}%
\paragraph{Pontos de dúvida}%
\begin{itemize}%
\item%
\textit{(statement)} Dá para ler "muito rápida" como irreversível. Nesse caso a curva P(V) de II e IV não é bem definida e não vale PV\^{}$\gamma$ = cte; é o que a resolução faz, com trechos tracejados. Também dá para ler como adiabática reversível, com curvas PV\^{}$\gamma$, que é o que o gabarito sugere com "adiabática". A conclusão de b não muda, porque $\Delta$U depende só de T. A resolução desenha os trechos II e IV mais inclinados que as isotermas, o que atende às duas leituras. Um corretor que espere adiabáticas desenhadas em linha cheia dificilmente descontaria.%
\item%
\textit{(solution)} Os dois usam W como trabalho sobre o gás ($\Delta$U = Q + W). Se o corretor usar trabalho realizado pelo gás ($\Delta$U = Q - W), os sinais de W1 a W4 se invertem, mas W2 + W4 = 0 e W\_tot = W1 + W3 continuam valendo. A resolução declara a convenção explicitamente.%
\end{itemize}%
\section*{Questão 3}%
\href{https://poliedroresolve.sistemapoliedro.com.br/vestibulares/ime/2027/ime-2-fase-dia-3-28-10-2026-2-fase-dia-3/dissertativa-3-quimica-2-fase-dia-3-ime-2027}{Resolução do gabarito}\par\smallskip%
\begin{longtable}{@{}p{0.06\textwidth} p{0.33\textwidth} p{0.33\textwidth} p{0.1\textwidth} p{0.08\textwidth}@{}}\toprule Item & Resolução & Gabarito & Veredito & Certo \\ \midrule
a & Ordem 2 em SO2 e ordem 0 em O3 (ordem global 2) & Ordem 2 em SO2, ordem 0 em O3 & igual & ambos \\
b & v = k{[}SO2{]}\^{}2 & v = k{[}SO2{]}\^{}2 & igual & ambos \\
c & k $\approx$ 1,89 L$\cdot$mol$^{-}$$^{1}$$\cdot$s$^{-}$$^{1}$ (a 300 K) & k = 1,888 L$\cdot$mol$^{-}$$^{1}$$\cdot$s$^{-}$$^{1}$ & equivalente & ambos \\
d & Ea = R ln(k2/k1)/(1/T1 - 1/T2) $\approx$ 4,93$\cdot$10$^3$ J/mol $\approx$ 4,9 kJ/mol & Ea = 594R $\approx$ 4,93 kJ/mol & igual & ambos \\
\bottomrule\end{longtable}%
\begin{itemize}%
\item%
\textbf{a)} Exp. 1$\to$2: {[}O3{]} cai à metade e v não muda, logo ordem 0 em O3. Exp. 2$\to$3: {[}SO2{]} triplica e v fica 9 vezes maior (1,062/0,118 = 9), logo ordem 2 em SO2.%
\item%
\textbf{b)} Sai direto das ordens do item a.%
\item%
\textbf{c)} 0,118/0,25$^2$ = 1,062/0,75$^2$ = 1,888 exato. A resolução arredonda para 1,89 e cita 1,888 na conferência. A diferença é só de arredondamento, e a unidade (ordem global 2) está correta nos dois.%
\item%
\textbf{d)} Refeito: k2/k1 = 7,56/1,888 $\approx$ 4,004 $\approx$ 4. Com ln 4 = 2$\cdot$0,693 = 1,386 e 1/300 - 1/1000 = 7/3000 K$^{-}$$^{1}$, Ea/R = 1,386$\cdot$3000/7 = 594 K. Com R = 8,3: Ea = 4930 J/mol $\approx$ 4,93 kJ/mol, o mesmo valor do gabarito (594R). Com 1,888 exato, Ea/R $\approx$ 594,6 K e Ea $\approx$ 4,935 kJ/mol, o que não muda nada. A resolução dá o valor final como 4,9 kJ/mol, mas também escreve explicitamente $\approx$ 4,93$\cdot$10$^3$ J/mol.%
\end{itemize}%
\paragraph{Pontos de dúvida}%
\begin{itemize}%
\item%
\textit{(statement)} O enunciado dá k(1000 K) = 7,56 'em módulo', sem unidade. Tanto a resolução quanto o gabarito supõem que a unidade é a mesma de k(300 K) (L$\cdot$mol$^{-}$$^{1}$$\cdot$s$^{-}$$^{1}$), ou seja, que a lei de velocidade não muda com a temperatura. É razoável, mas não está dito.%
\item%
\textit{(statement)} A reação ocorre sob luz UV (fotoquímica). A rigor, aplicar Arrhenius a uma reação fotoiniciada é discutível. Tanto a resolução quanto o gabarito seguem o que a questão pede e usam Arrhenius diretamente. A resolução comenta que uma Ea baixa combina com a iniciação por UV.%
\item%
\textit{(solution)} Resposta final do item d em 4,9 kJ/mol (2 algarismos significativos), enquanto o gabarito dá 4,93 kJ/mol. A resolução também mostra 4,93$\cdot$10$^3$ J/mol, então a diferença é só de arredondamento.%
\item%
\textit{(statement)} Supõe{-}se que a 'velocidade inicial' da tabela é a velocidade da reação. Como a estequiometria é toda 1:1, ela coincide com a velocidade de consumo de qualquer reagente, e por isso não há fator a corrigir em k.%
\end{itemize}%
\section*{Questão 4}%
\href{https://poliedroresolve.sistemapoliedro.com.br/vestibulares/ime/2027/ime-2-fase-dia-3-28-10-2026-2-fase-dia-3/dissertativa-4-quimica-2-fase-dia-3-ime-2027}{Resolução do gabarito}\par\smallskip%
\begin{longtable}{@{}p{0.06\textwidth} p{0.33\textwidth} p{0.33\textwidth} p{0.1\textwidth} p{0.08\textwidth}@{}}\toprule Item & Resolução & Gabarito & Veredito & Certo \\ \midrule
a & Ligações (pontes) de hidrogênio O--H$\cdot$$\cdot$$\cdot$O=C, duas por dímero, fechando um anel de 8 átomos. & Duas ligações de hidrogênio entre as carboxilas. & igual & ambos \\
b & Dímero cíclico (48 elétrons de valência) com duas pontes de hidrogênio em pontilhado (figura TikZ). & Dímero cíclico com duas ligações de H (estrutura em imagem). & equivalente & ambos \\
c & Kc $\approx$ 1,6$\cdot$10$^2$ L/mol (164 L/mol), com {[}M{]} = 0,016 mol/L e {[}D{]} = 0,042 mol/L, usando c0 = 0,10 mol/L do item d. Também apresenta a forma geral Kc = 16,4/c0. & Kc = 0,42V/0,16$^2$ = 16,4$\cdot$V (V em L). Sem valor numérico, porque o enunciado não dá volume nem concentração. & diferente & indefinido \\
d & Tf $\approx$ 5,11 $^\circ$C (i = 0,58; W de partículas = 0,0671 mol/kg; $\Delta$T $\approx$ 0,34 K; benzeno $\approx$ 0,865 kg por litro de solução). & i = 0,58; W = 0,1156 mol/kg; $\Delta$Tc = 0,343 $^\circ$C; T = 5,11 $^\circ$C. & igual & ambos \\
\bottomrule\end{longtable}%
\begin{itemize}%
\item%
\textbf{a)} Mesma resposta: duas ligações de hidrogênio entre o H da hidroxila de uma molécula e o O carbonílico da outra.%
\item%
\textbf{b)} Conferi a figura da resolução. Cada C liga{-}se a um CH3, a um O por ligação dupla e a um O por ligação simples. Cada O--H forma ponte de H (pontilhada) com o O carbonílico da outra molécula, fechando o anel de 8 átomos C=O$\cdot$$\cdot$$\cdot$H--O--C=O$\cdot$$\cdot$$\cdot$H--O. Há 2 pares isolados em cada O, e a contagem de 24 e$^{-}$ por monômero (8 pares ligantes + 4 isolados) está certa. A única ressalva é que o metila aparece condensado (H3C), sem as ligações C--H desenhadas. A imagem do gabarito não está disponível, mas a descrição bate.%
\item%
\textbf{c)} Refiz a conta: n\_D = 0,42 e n\_M = 1 - 2(0,42) = 0,16 por mol de ácido, então Kc = (0,42/V)/(0,16/V)$^2$ = 16,406$\cdot$V = 16,4/c0. A expressão geral da resolução (16,4/c0) é idêntica à do gabarito, porque V = 1/c0 por mol de ácido. Com c0 = 0,10 mol/L (V = 10 L), sai 164,06 L/mol, que é o valor da resolução. A divergência é só sobre a hipótese que fixa a concentração. A resolução argumenta que, com o mesmo Kc, manter {[}D{]}/{[}M{]} = Kc$\cdot${[}M{]} obriga a manter {[}M{]} e, portanto, c0 = 0,10 mol/L. O argumento é válido apenas à mesma temperatura, e o enunciado não diz que o equilíbrio de referência está na T da solução do item d ($\approx$ 5 $^\circ$C). O gabarito é o mais rigoroso com os dados fornecidos. Já o enunciado pede 'o valor', sugerindo um número, que pode ser 164 (c0 = 0,10) ou até 16,4 (base 1 L). Não dá para decidir sem o critério oficial.%
\item%
\textbf{d)} Refiz a conta. O ácido ocupa 0,10$\cdot$60,05/1050 = 5,72$\cdot$10$^{-}$$^3$ L, então a massa de benzeno é 0,9943$\cdot$0,87 = 0,865 kg e W = 0,10/0,865 = 0,1156 mol/kg. Com i = 1 - 0,84/2 = 0,58, $\Delta$T = 5,12$\cdot$0,1156$\cdot$0,58 = 0,3433 K e Tf = 5,45 - 0,343 = 5,107 $\approx$ 5,11 $^\circ$C. O gabarito escreve a molalidade estequiométrica (0,1156) e multiplica por i. A resolução usa direto a molalidade de partículas (0,058/0,865 = 0,0671), o que dá o mesmo resultado.%
\end{itemize}%
\paragraph{Pontos de dúvida}%
\begin{itemize}%
\item%
\textit{(statement)} Item c: o enunciado não dá a concentração (ou o volume) do equilíbrio de referência, e Kc = 16,4/c0 depende dela. O gabarito Poliedro deixa a resposta como Kc = 16,4$\cdot$V. A resolução fixa c0 = 0,10 mol/L supondo que a solução do item d é o próprio equilíbrio, e obtém 164 L/mol. Para isso precisa supor a mesma temperatura, o que o enunciado não garante: a solução do item d está no ponto de fusão, $\approx$ 5 $^\circ$C, e a T do equilíbrio original não é dada. O gabarito oficial do IME pode esperar 164 (c0 = 0,10), 16,4 (base 1 L) ou a expressão literal. Vale levar a um professor que critério de correção será aceito.%
\item%
\textit{(solution)} A resolução chega ao valor numérico de Kc pela hipótese implícita de mesma temperatura (Kc constante). Ela escreve 'mesma T', mas não discute que a solução do item d está a $\approx$ 5 $^\circ$C. O argumento é coerente, mas é uma hipótese extra, não um dado do problema.%
\item%
\textit{(statement)} A frase 'cada mol de monômero gera 0,42 mol de dímero' admite outra leitura: {[}D{]}/{[}M{]} = 0,42 no equilíbrio, o que daria Kc $\approx$ 7,7 L/mol e Tf $\approx$ 4,99 $^\circ$C. Resolução e gabarito adotam a mesma leitura (n\_D = 0,42 por mol de ácido, i = 0,58), então não há divergência. A resolução descarta a alternativa em parte com um dado de literatura (Kc \textasciitilde{}10$^2$ L/mol), que é externo à prova.%
\item%
\textit{(solution)} Item b: na estrutura de Lewis os grupos metila aparecem condensados (H3C/CH3), sem as ligações C--H explícitas. Uma correção rígida de 'estrutura de Lewis' poderia exigir todos os átomos e ligações.%
\item%
\textit{(key)} Item d: a massa de benzeno ($\approx$ 0,865 kg por litro de solução) supõe volumes aditivos, tanto na resolução quanto no gabarito (W = 0,1156 mol/kg). Aproximar para 0,87 kg dá o mesmo Tf = 5,11 $^\circ$C. O gabarito expressa $\Delta$T em $^\circ$C, o que é equivalente a K para uma variação.%
\end{itemize}%
\section*{Questão 5}%
\href{https://poliedroresolve.sistemapoliedro.com.br/vestibulares/ime/2027/ime-2-fase-dia-3-28-10-2026-2-fase-dia-3/dissertativa-5-quimica-2-fase-dia-3-ime-2027}{Resolução do gabarito}\par\smallskip%
\begin{longtable}{@{}p{0.06\textwidth} p{0.33\textwidth} p{0.33\textwidth} p{0.1\textwidth} p{0.08\textwidth}@{}}\toprule Item & Resolução & Gabarito & Veredito & Certo \\ \midrule
unico & n{-}heptano (C7H16): $\Delta$n\_gás = C($\Delta$T\_V - $\Delta$T\_P)/(RT) = 450$\cdot$16,6/(8,3$\cdot$300) = +3, e $\Delta$n = y/4 - 1 = 3 só para y = 16 & $\Delta$n\_gás = 450$\cdot$16,6/(8,3$\cdot$300) = 3 $\Rightarrow$ n{-}heptano (C7H16 + 11 O2 $\to$ 7 CO2 + 8 H2O) & igual & ambos \\
\bottomrule\end{longtable}%
\begin{itemize}%
\item%
\textbf{unico)} A conclusão é a mesma: n{-}heptano. Refiz as contas: 450$\cdot$16,6 = 7470 e 8,3$\cdot$300 = 2490, o que dá $\Delta$n = 3,00. Para C7H16 + 11 O2 $\to$ 7 CO2 + 8 H2O(g), há 15 mol de produtos gasosos contra 12 de reagentes, logo $\Delta$n = +3. O sinal também confere. Como q\_V = $\Delta$U e q\_P = $\Delta$H, temos C($\Delta$T\_V - $\Delta$T\_P) = $\Delta$H - $\Delta$U = $\Delta$n$\cdot$RT. Com $\Delta$n > 0, a liberação de calor a volume constante é maior, então $\Delta$T\_V > $\Delta$T\_P, como diz o item (v). Para os outros candidatos, a fórmula geral $\Delta$n = y/4 - 1 dá: acetileno -0,5 (-2,77 K), p{-}xileno +1,5 (+8,3 K) e propeno +0,5 (+2,77 K). Nenhum chega a 16,6 K. A transcrição do enunciado confere com a página 3 da prova.%
\end{itemize}%
\paragraph{Pontos de dúvida}%
\begin{itemize}%
\item%
\textit{(statement)} O enunciado é idealizado. O item (i) manda considerar a reação a 27 $^\circ$C (T = 300 K em $\Delta$(pV) = $\Delta$n$\cdot$RT), mas os calorímetros esquentam muito. Com C = 450 J/K e cerca de 4,5$\cdot$10\^{}6 J/mol, o $\Delta$T absoluto seria da ordem de 10\^{}4 K. A resolução comenta isso, e só a diferença $\Delta$T\_V - $\Delta$T\_P entra na conta.%
\item%
\textit{(statement)} O p{-}xileno e o n{-}heptano são líquidos a 27 $^\circ$C e 1 atm, mas o enunciado manda tratar tudo como gás (item ii), com H2O em vapor. A resolução adota isso de forma explícita, e é o que leva a $\Delta$n = y/4 - 1. Se a água fosse líquida, $\Delta$n mudaria e nenhum candidato daria +3.%
\item%
\textit{(solution)} No comentário de ordem de grandeza, a resolução usa 4,8$\cdot$10\^{}6 J para o calor de combustão do heptano. Esse valor corresponde mais ao poder calorífico com água líquida; com vapor, seria cerca de 4,5$\cdot$10\^{}6 J. Isso não afeta a resposta, já que o comentário é só ilustrativo.%
\end{itemize}%
\section*{Questão 6}%
\href{https://poliedroresolve.sistemapoliedro.com.br/vestibulares/ime/2027/ime-2-fase-dia-3-28-10-2026-2-fase-dia-3/dissertativa-6-quimica-2-fase-dia-3-ime-2027}{Resolução do gabarito}\par\smallskip%
\begin{longtable}{@{}p{0.06\textwidth} p{0.33\textwidth} p{0.33\textwidth} p{0.1\textwidth} p{0.08\textwidth}@{}}\toprule Item & Resolução & Gabarito & Veredito & Certo \\ \midrule
a & CH3COOH(l) + C2H5OH(l) ? CH3COOC2H5(l) + H2O(l), proporção 1:1:1:1 & CH3COOH(l) + C2H5OH(l) ? CH3COOC2H5(l) + H2O(l) & igual & ambos \\
b & $\alpha$(ácido acético) = 1,6/4,0 = 0,40 (40\%) & $\alpha$(ácido acético) = 1,6/4 = 0,4 (40\%) & igual & ambos \\
c & Kc = (1,6$\cdot$1,6)/(2,4$\cdot$0,4) = 8/3 $\approx$ 2,67 (adimensional, V se cancela) & Kc = (1,6$\cdot$1,6)/(2,4$\cdot$0,4) = 8/3 & igual & ambos \\
\bottomrule\end{longtable}%
\begin{itemize}%
\item%
\textbf{a)} Esterificação de Fischer escrita do mesmo jeito, com os mesmos estados físicos. Conferi o balanço: C4H10O3 dos dois lados.%
\item%
\textbf{b)} n(ácido) = 240/60 = 4,0 mol, e o ácido consumido é 1,6 mol (estequiometria 1:1), logo $\alpha$ = 0,40. Com as massas molares da tabela (60,05 g/mol) dá $\alpha$ $\approx$ 0,4003, sem diferença que importe. A resolução ainda observa, à parte, que a conversão do etanol seria 80\%, o que não atrapalha a resposta.%
\item%
\textbf{c)} No equilíbrio: ácido 2,4 mol, etanol 0,4 mol, éster 1,6 mol, água 1,6 mol. Como $\Delta$n = 0, o volume se cancela e Kc = 2,56/0,96 = 8/3 $\approx$ 2,67. Conferi em python: as raízes de 5x$^2$ - 48x + 64 = 0 são 1,6 e 8, o que confirma a coerência. Com massas molares exatas (60,05 e 46,07) o resultado seria Kc $\approx$ 2,69, diferença que é só de arredondamento.%
\end{itemize}%
\paragraph{Pontos de dúvida}%
\begin{itemize}%
\item%
\textit{(solution)} A água entra na expressão de Kc, pois o meio é uma única fase líquida e não aquosa. A resolução declara essa hipótese, e o gabarito faz o mesmo. Se alguém tirasse a água por tratá{-}la como 'líquido puro', obteria outro valor; vale confirmar com o professor que a convenção de incluí{-}la é a esperada pelo IME, como é o usual na esterificação.%
\item%
\textit{(statement)} O valor de Kc depende de como se arredondam as massas molares: com 60 e 46 g/mol sai 8/3 $\approx$ 2,67; com os valores da tabela (60,05 e 46,07 g/mol) sai $\approx$ 2,69. A prova está claramente pensada para 4,0 e 2,0 mol, então a diferença não deve contar na nota.%
\end{itemize}%
\section*{Questão 7}%
\href{https://poliedroresolve.sistemapoliedro.com.br/vestibulares/ime/2027/ime-2-fase-dia-3-28-10-2026-2-fase-dia-3/dissertativa-7-quimica-2-fase-dia-3-ime-2027}{Resolução do gabarito}\par\smallskip%
\begin{longtable}{@{}p{0.06\textwidth} p{0.33\textwidth} p{0.33\textwidth} p{0.1\textwidth} p{0.08\textwidth}@{}}\toprule Item & Resolução & Gabarito & Veredito & Certo \\ \midrule
a & $\alpha$ = 1, $\beta$ = 2, $\gamma$ = -1 (adota $\nu$$_{0}$ do exp. 4 = 1,28$\times$10$^{-}$$^3$, supondo erro de expoente na impressão) & $\alpha$ = 1, $\beta$ = 2, $\gamma$ = -1 (supondo erro de digitação: 1,28$\cdot$10$^{-}$$^3$ em vez de 1,28$\cdot$10$^{-}$$^{6}$) & igual & ambos \\
b & k = 1,25$\times$10$^{-}$$^2$ L$\cdot$mol$^{-}$$^{1}$$\cdot$s$^{-}$$^{1}$ & k = 1,25$\cdot$10$^{-}$$^2$ L$\cdot$mol$^{-}$$^{1}$$\cdot$s$^{-}$$^{1}$ & igual & ambos \\
c & $\nu$$_{0}$ = 1,25$\times$10$^{-}$$^2$ {[}A{]}{[}B{]}$^2$/{[}C{]} (concentrações em mol/L, $\nu$$_{0}$ em mol L$^{-}$$^{1}$ s$^{-}$$^{1}$) & v = k{[}A{]}{[}B{]}$^2${[}C{]}$^{-}$$^{1}$ & equivalente & ambos \\
d & ordem global = 1 + 2 - 1 = 2 & ordem global = 2 & igual & ambos \\
\bottomrule\end{longtable}%
\begin{itemize}%
\item%
\textbf{a)} Conferi o enunciado no PDF (p. 4): a tabela traz mesmo $\nu$$_{4}$ = 1,28$\times$10$^{-}$$^{6}$. As razões em relação ao exp. 1 dão $\alpha$+2$\beta$+3$\gamma$=2 e 2$\alpha$+$\beta$+4$\gamma$=0. O exp. 4 dá 3$\alpha$+4$\beta$+$\gamma$=log$_{2}$($\nu$$_{4}$/$\nu$$_{1}$). O determinante é 20 (conferido). Com $\nu$$_{4}$ = 1,28$\times$10$^{-}$$^3$ a razão é 1024 = 2$^{1}$$^{0}$ e a solução é exatamente (1, 2, -1). Com o valor impresso sai (-1,49; 1,00; 0,49), que não tem sentido físico. A resolução e o gabarito fazem a mesma correção. A resolução ainda confirma o resultado pelo estado estacionário: no limite k$_{3}${[}C{]} ? k$_{2}${[}B{]}, as razões 4, 1 e 1024 batem com os experimentos 2 e 3 e com o exp. 4 corrigido.%
\item%
\textbf{b)} Pelo exp. 1, k = $\nu$$_{0}${[}C{]}/({[}A{]}{[}B{]}$^2$) = (1,25$\times$10$^{-}$$^{6}$$\cdot$10$^{-}$$^2$)/(10$^{-}$$^2$$\cdot$10$^{-}$$^{4}$) = 1,25$\times$10$^{-}$$^2$. A unidade sai (mol L$^{-}$$^{1}$ s$^{-}$$^{1}$)(mol L$^{-}$$^{1}$)/(mol L$^{-}$$^{1}$)$^3$ = L mol$^{-}$$^{1}$ s$^{-}$$^{1}$. Usando esse k, os exp. 2 e 3 dão 5,00$\times$10$^{-}$$^{6}$ e 1,25$\times$10$^{-}$$^{6}$ (iguais à tabela) e o exp. 4 dá 1,28$\times$10$^{-}$$^3$ (igual ao valor corrigido).%
\item%
\textbf{c)} É a mesma lei. A resolução apenas substitui o valor numérico de k, e {[}C{]}$^{-}$$^{1}$ equivale a 1/{[}C{]}.%
\item%
\textbf{d)} $\alpha$+$\beta$+$\gamma$ = 2, coerente com a unidade L mol$^{-}$$^{1}$ s$^{-}$$^{1}$ de k.%
\end{itemize}%
\paragraph{Pontos de dúvida}%
\begin{itemize}%
\item%
\textit{(statement)} Pelo PDF original, a tabela impressa traz $\nu$$_{0}$ = 1,28$\times$10$^{-}$$^{6}$ mol L$^{-}$$^{1}$ s$^{-}$$^{1}$ no exp. 4. Com esse valor os expoentes saem não inteiros ($\alpha$ $\approx$ -1,49; $\beta$ $\approx$ 1,00; $\gamma$ $\approx$ 0,49). Só com 1,28$\times$10$^{-}$$^3$ (razão 2$^{1}$$^{0}$ = 1024) se chega a (1, 2, -1). A resolução e o gabarito fazem a mesma correção. Convém o professor verificar se a banca publicou errata ou se vai aceitar outras abordagens: usar o valor literal, ou usar só os exp. 1 a 3 com o mecanismo, que já bastam para chegar a (1, 2, -1) se a ordem em A for 1.%
\item%
\textit{(solution)} A resolução justifica o expoente negativo por inibição: C compete por X na etapa III, e $\gamma$ = -1 vem do estado estacionário no limite k$_{3}${[}C{]} ? k$_{2}${[}B{]}. Isso pressupõe k$_{3}$/k$_{2}$ ? 8 também no exp. 4, em que {[}B{]}/{[}C{]} = 8. É uma hipótese implícita razoável, mas não é dita explicitamente.%
\item%
\textit{(solution)} Há um deslize de redação na eliminação: '2$\cdot$(2/1 da 2ª)-(1ª)'. A conta está certa: 2$\cdot$(2ª eq.) - (1ª eq.) dá 3$\alpha$ + 5$\gamma$ = -2.%
\end{itemize}%
\section*{Questão 8}%
\href{https://poliedroresolve.sistemapoliedro.com.br/vestibulares/ime/2027/ime-2-fase-dia-3-28-10-2026-2-fase-dia-3/dissertativa-8-quimica-2-fase-dia-3-ime-2027}{Resolução do gabarito}\par\smallskip%
\begin{longtable}{@{}p{0.06\textwidth} p{0.33\textwidth} p{0.33\textwidth} p{0.1\textwidth} p{0.08\textwidth}@{}}\toprule Item & Resolução & Gabarito & Veredito & Certo \\ \midrule
a & Y é o meso{-}3,4{-}dimetil{-}hexano, (3R,4S) $\equiv$ (3S,4R). Foi desenhado em zigue{-}zague anti, com cunha e tracejado (CH3 em cunha em C3 e tracejado em C4, e centro de inversão entre C3 e C4), e em projeção de Fischer, com os dois CH3 do mesmo lado. & Y é a forma meso do 3,4{-}dimetil{-}hexano (a estrutura 3D está em imagem). & igual & ambos \\
b & X é o 2{-}clorobutano, CH3CH2CHClCH3, que precisa ser racêmico ou pelo menos fornecer um grupo R e um S para gerar a forma meso. & X é o 2{-}clorobutano. & igual & ambos \\
\bottomrule\end{longtable}%
\begin{itemize}%
\item%
\textbf{a)} Pela reação de Wurtz, 2 R--Cl + 2 Na $\to$ R--R + 2 NaCl. Então Y é simétrico e tem um número par de carbonos. Com 2, 4 ou 6 C (etano, butano, hexano ou 2,3{-}dimetilbutano) não há carbono quiral. Com 8 C, só o sec{-}butila gera dois estereocentros (C3 e C4). Como Y não desvia a luz polarizada, ele é a forma meso. Conferi as configurações CIP com coordenadas 3D tiradas do próprio desenho e calculei o produto misto. As prioridades são C4 > C2 > CH3 > H em C3, e C3 > C5 > CH3 > H em C4. O zigue{-}zague dá C3 = R e C4 = S, e a projeção de Fischer também dá (3R,4S). Pelas coordenadas, o confôrmero anti desenhado é centrossimétrico: a inversão pelo ponto médio de C3--C4 leva CH3 em CH3 e H em H. Portanto o desenho representa de fato a forma meso. Não dá para comparar com a imagem do gabarito, mas a estrutura é a mesma.%
\item%
\textbf{b)} Como R = sec{-}butila, X = CH3CH2CHClCH3. A cadeia principal tem 4 C e o Cl recebe o menor localizador, 2. O nome IUPAC é 2{-}clorobutano, idêntico ao do gabarito. O comentário sobre X ser racêmico é um complemento correto e não altera a resposta.%
\end{itemize}%
\paragraph{Pontos de dúvida}%
\begin{itemize}%
\item%
\textit{(solution)} Na prática, a reação de Wurtz com haletos secundários como o 2{-}clorobutano rende pouco: competem eliminação e desproporcionamento, que geram butenos e butano. A resolução trata a reação como acoplamento limpo, que é a hipótese pedida pelo enunciado. Também não comenta que a forma meso sai misturada com o par (3R,4R)/(3S,4S), que ela própria cita, e que precisaria ser separada.%
\item%
\textit{(statement)} O enunciado chama C3 e C4 de "carbonos quirais" num composto meso, que é aquiral. Isso segue o uso escolar comum de "carbono assimétrico". A rigor, são estereocentros (carbonos assimétricos) numa molécula aquiral. Não muda a resposta.%
\item%
\textit{(key)} O gabarito traz a estrutura 3D só como imagem, que não foi fornecida. Por isso não é possível conferir se o desenho do gabarito coincide com a conformação ou projeção escolhida pela resolução. Qualquer representação 3D correta da forma meso deveria valer.%
\end{itemize}%
\section*{Questão 9}%
\href{https://poliedroresolve.sistemapoliedro.com.br/vestibulares/ime/2027/ime-2-fase-dia-3-28-10-2026-2-fase-dia-3/dissertativa-9-quimica-2-fase-dia-3-ime-2027}{Resolução do gabarito}\par\smallskip%
\begin{longtable}{@{}p{0.06\textwidth} p{0.33\textwidth} p{0.33\textwidth} p{0.1\textwidth} p{0.08\textwidth}@{}}\toprule Item & Resolução & Gabarito & Veredito & Certo \\ \midrule
unico & M = 106 g/mol; A = p{-}xileno (1,4{-}dimetilbenzeno), C8H10 & M = 106 g/mol; A = p{-}xileno (1,4{-}dimetilbenzeno) & igual & ambos \\
\bottomrule\end{longtable}%
\begin{itemize}%
\item%
\textbf{unico)} Refiz a conta: W = $\Delta$T/Kc = 1,28/5,12 = 0,25 mol/kg; n = 0,25 $\times$ 2 = 0,50 mol; M = 53/0,50 = 106 g/mol. Com essa massa molar, as opções aromáticas plausíveis são C8H10 (106,17 g/mol: etilbenzeno e o{-}, m{-}, p{-}xileno) e C7H6O (106,12 g/mol: benzaldeído). O Tollens negativo exclui o benzaldeído. Na halogenação, só o p{-}xileno dá um único monoderivado no anel (os 4 H aromáticos são equivalentes) e também um único na cadeia lateral. A conferência inversa dá $\Delta$T = 5,12$\cdot$(53/106,17)/2 $\approx$ 1,278 $^\circ$C. A resolução e o gabarito chegam à mesma estrutura e à mesma massa molar.%
\end{itemize}%
\paragraph{Pontos de dúvida}%
\begin{itemize}%
\item%
\textit{(statement)} O enunciado diz "um único produto possível" sem dizer se a halogenação é no anel (X2/FeX3) ou na cadeia lateral (X2/luz), nem se conta só a monossubstituição. Na leitura radicalar, o o{-}xileno e o m{-}xileno também dariam um único monoderivado. A resolução trata isso com cuidado: o p{-}xileno é o único que atende às duas leituras. E se contarem também di{-} ou tri{-}halogenações, nem o p{-}xileno dá um produto só. A leitura implícita, monoderivado, é a usual.%
\item%
\textit{(solution)} A busca de fórmulas assume um anel benzênico com substituintes de C, H e O. Não considera heteroaromáticos com N (por exemplo, derivados de piridina com M = 106) nem outros heteroátomos. Isso é razoável no nível do IME e não muda a resposta esperada.%
\item%
\textit{(solution)} A resolução declara hipóteses não dadas no enunciado: solução diluída, soluto não volátil e i = 1 (A não se associa nem se dissocia no benzeno). Elas são razoáveis e foram justificadas (x\_A $\approx$ 0,019).%
\end{itemize}%
\section*{Questão 10}%
\href{https://poliedroresolve.sistemapoliedro.com.br/vestibulares/ime/2027/ime-2-fase-dia-3-28-10-2026-2-fase-dia-3/dissertativa-10-quimica-2-fase-dia-3-ime-2027}{Resolução do gabarito}\par\smallskip%
\begin{longtable}{@{}p{0.06\textwidth} p{0.33\textwidth} p{0.33\textwidth} p{0.1\textwidth} p{0.08\textwidth}@{}}\toprule Item & Resolução & Gabarito & Veredito & Certo \\ \midrule
unico & t = t$\frac12$$\cdot$ln(1,024)/ln 2 $\approx$ 4,9$\times$10$^{8}$ anos (com ln 1,024 = 7 ln2 - 3 ln5 $\approx$ 0,024; t/t$\frac12$ $\approx$ 0,0346) & t/t$\frac12$ = 0,03463 $\Rightarrow$ t $\approx$ 4,88$\cdot$10$^{8}$ anos & equivalente & ambos \\
\bottomrule\end{longtable}%
\begin{itemize}%
\item%
\textbf{unico)} Os dois partem de N\_Pb/N\_Th = e\^{}\{$\lambda$t\} - 1 = 0,024 e tomam ln(1,024) $\approx$ 0,024 com os dados fornecidos (7$\cdot$0,693 - 3$\cdot$1,609 = 0,024). Assim t/t$\frac12$ = 0,024/0,693 = 0,03463 e t = 1,41$\times$10$^{1}$$^{0}$ $\times$ 0,03463 = 4,883$\times$10$^{8}$ anos, que é o número do gabarito. A resolução arredonda para 4,9$\times$10$^{8}$ anos. A diferença é só de algarismos significativos. Conferi em Python: com o ln exato (0,023717), dá t = 4,82$\times$10$^{8}$ anos, valor que a própria resolução cita como \textasciitilde{}4,8$\times$10$^{8}$. Os dois estão certos dentro da precisão dos dados. Confirmei no PDF (p. 5) que o enunciado bate com a transcrição.%
\end{itemize}%
\paragraph{Pontos de dúvida}%
\begin{itemize}%
\item%
\textit{(solution)} O resultado oscila entre a aproximação ln(1,024) $\approx$ 0,024, que dá 4,88$\times$10$^{8}$ anos, e o ln exato, que dá 4,82$\times$10$^{8}$ anos. A resolução declara 4,9$\times$10$^{8}$ anos e cita 4,8$\times$10$^{8}$ anos como conferência. As duas formas são aceitáveis, e a banca deve aceitar qualquer valor na faixa de \textasciitilde{}4,8--4,9$\times$10$^{8}$ anos.%
\item%
\textit{(solution)} Imprecisão de redação na conferência: 'apenas 2,4\% do Th ter virado Pb'. A fração do Th inicial que decaiu é 0,024/1,024 $\approx$ 2,34\%. O 0,024 é a razão Pb/Th atual, não a fração decaída. Isso não afeta a conta.%
\item%
\textit{(statement)} Os teores de ThO$_{2}$ (9\%) e de U$_{3}$O$_{8}$ (0,35\%) e a meia{-}vida do $^2$$^3$$^{8}$U são distratores: só a razão $^2$$^{0}$$^{8}$Pb/$^2$$^3$$^2$Th é necessária para a idade. A resolução reconhece isso de forma explícita. Fica implícita a hipótese de equilíbrio secular na série do Th (1 átomo de Th $\to$ 1 átomo de $^2$$^{0}$$^{8}$Pb), que a resolução declara.%
\end{itemize}%
\end{document}